\documentclass[11pt]{article}
\usepackage[a4paper,margin=0.65in]{geometry}
\usepackage[T1]{fontenc}
\usepackage{amsmath,amssymb,amsfonts}
\usepackage{graphicx}
\usepackage{pdfpages}
\usepackage[english,shorthands=off]{babel}
\usepackage{fancyhdr}
\usepackage[bookmarks=false,colorlinks=false]{hyperref}
\usepackage[protrusion=true,expansion=false]{microtype}
\emergencystretch=3em
\tolerance=600
\providecommand{\modd}{\mathrm{mod}}
\providecommand{\Disc}{\mathrm{Disc}}
\providecommand{\canont}{}
\providecommand{\asterism}{*}
\providecommand{\Hessian}{H}
\providecommand{\tome}{}
\providecommand{\page}{p.}
\pagestyle{fancy}\fancyhf{}\fancyhead[L]{Pages 476--500}\fancyhead[R]{Cayley --- Collected Papers, Vol. X}\renewcommand{\headrulewidth}{0.4pt}
\setlength{\headheight}{15pt}

\begin{document}

\begin{center}
\textbf{\large Cayley, Collected Mathematical Papers, Vol. X}\\[2pt]
\textit{Pages 476--500 (scan); papers 703 (continued) and 704}
\end{center}

\bigskip
\hrule
\bigskip

\noindent\textbf{[Continuation of Paper 703: On the Addition of the Double $\vartheta$-Functions.]}

\bigskip

I attach the numbers $1,2,3,4,5,6$ to the variables $x,y,z,w,p,q$, respectively: and write
\[
A_{12} = \sqrt{a-x\mathbin{.}a-y},\quad A_{34} = \sqrt{a-z\mathbin{.}a-w},\quad A_{56} = \sqrt{a-p\mathbin{.}a-q};
\]
(six equations),
\[
AB_{12} = \tfrac{1}{x-y}\bigl\{\sqrt{a-x\mathbin{.}b-x\mathbin{.}f-x\mathbin{.}c-y\mathbin{.}d-y\mathbin{.}e-y} - \sqrt{a-y\mathbin{.}b-y\mathbin{.}f-y\mathbin{.}c-x\mathbin{.}d-x\mathbin{.}e-x}\bigr\};\ \text{etc.}
\]
(ten equations),

where it is to be borne in mind that $AB$ is an abbreviation for $ABF\mathbin{.}CDE$, and so in other cases, the letter $F$ belonging always to the expressed duad: there are thus in all the sixteen functions $A,B,C,D,E,F,AB,AC,AD,AE,BC,BD,BE,CD,CE,DE$, these being functions of $x$ and $y$, of $z$ and $w$, and of $p$ and $q$, according as the suffix is $12$, $34$, or $56$.

It is to be shown that the $16$ functions $A_{56}$, $AB_{56}$ of $p$ and $q$ can be by means of the given equations expressed as proportional to rational and integral functions of the $16$ functions $A_{12}$, $AB_{12}$, $A_{34}$, $AB_{34}$ of $x$ and $y$, and of $z$ and $w$ respectively: and it is clear that in so expressing them we have in effect the solution of the problem of the addition of the double $\vartheta$-functions.

I use when convenient the abbreviated notations
\[
a-x = a_1,\ a-y = a_2,\quad a-z = a_3,\ a-w = a_4,\quad \text{etc.,}
\]
we have of course
\[
\theta_{12} = x-y,\ \theta_{34} = z-w,\ \theta_{56} = p-q;
\]
\[
X = a_1 b_1 c_1 d_1 e_1 f_1,
\]
\[
\sqrt{ab} = \tfrac{1}{x-y}\bigl\{\sqrt{a_1 b_1 f_1 c_2 d_2 e_2} - \sqrt{a_2 b_2 f_2 c_1 d_1 e_1}\bigr\},\ \text{etc.}
\]

Proceeding to the investigation, the equations between the variables are obviously those obtained by the elimination of the arbitrary multipliers $\alpha,\beta,\gamma,\delta,\varepsilon$ from the six equations obtained from
\[
\alpha a^2 + \beta a^2 + \gamma a + \delta = \varepsilon\sqrt{a},
\]
by writing therein for $a$ the values $x,y,z,w,p,q$ successively; we may consider the four equations
\begin{align*}
\alpha x^3 + \beta x^2 + \gamma x + \delta &= \varepsilon\sqrt{X},\\
\alpha y^3 + \beta y^2 + \gamma y + \delta &= \varepsilon\sqrt{Y},\\
\alpha z^3 + \beta z^2 + \gamma z + \delta &= \varepsilon\sqrt{Z},\\
\alpha w^3 + \beta w^2 + \gamma w + \delta &= \varepsilon\sqrt{W},
\end{align*}
as serving to determine the ratios $\alpha:\beta:\gamma:\delta:\varepsilon$ in terms of $x,y,z,w$; and we have then for the determination of $p,q$ the remaining two equations
\[
\alpha p^3 + \beta p^2 + \gamma p + \delta = \varepsilon\sqrt{P},
\]
which two equations may be replaced by the identity
\[
(\alpha\theta^3 + \beta\theta^2 + \gamma\theta + \delta)^2 - \varepsilon^2 \theta = \alpha^2(\theta - x)(\theta - y)(\theta - z)(\theta - w)(\theta - p)(\theta - q).
\]
Writing herein $\theta$ any one of the values $a,b,c,d,e,f$, for instance $\theta = a$, and taking the square root of each side, we have
\[
\alpha a^3 + \beta a^2 + \gamma a + \delta = \sqrt{a^6 - \varepsilon^2}\sqrt{a-x\mathbin{.}a-y}\sqrt{a-z\mathbin{.}a-w}\sqrt{a-p\mathbin{.}a-q},
\]
or as this may be written
\[
\alpha a^3 + \beta a^2 + \gamma a + \delta = \sqrt{a^6 - \varepsilon^2}\, A_{12}\, A_{34}\, A_{56},
\]
which equation when reduced gives the proportional value of $A_{56}$.

For the reduction we require the value of $\alpha a^3 + \beta a^2 + \gamma a + \delta$: calling this for the moment $\Omega$, we join to the four equations a fifth equation
\[
\alpha a^3 + \beta a^2 + \gamma a + \delta = \Omega.
\]
Eliminating $\alpha,\beta,\gamma,\delta$, we find
\[
\begin{vmatrix} \varepsilon\sqrt{X} & x^3 & x^2 & x & 1 \\ \varepsilon\sqrt{Y} & y^3 & y^2 & y & 1 \\ \varepsilon\sqrt{Z} & z^3 & z^2 & z & 1 \\ \varepsilon\sqrt{W} & w^3 & w^2 & w & 1 \\ \Omega & a^3 & a^2 & a & 1 \end{vmatrix} = 0,
\]
or, what is the same thing,
\[
\Omega\begin{vmatrix} x^3 & x^2 & x & 1 \\ y^3 & y^2 & y & 1 \\ z^3 & z^2 & z & 1 \\ w^3 & w^2 & w & 1 \end{vmatrix} + \varepsilon \begin{vmatrix} \sqrt{X} & x^3 & x^2 & x & 1 \\ \sqrt{Y} & y^3 & y^2 & y & 1 \\ \sqrt{Z} & z^3 & z^2 & z & 1 \\ \sqrt{W} & w^3 & w^2 & w & 1 \\ & a^3 & a^2 & a & 1 \end{vmatrix} = 0;
\]
viz.\ this is
\begin{align*}
\Omega\mathbin{.}\overline{x-y}\mathbin{.}\overline{x-z}\mathbin{.}\overline{x-w}\mathbin{.}\overline{y-z}\mathbin{.}\overline{y-w}\mathbin{.}\overline{z-w} =\ &-\varepsilon\bigl\{\sqrt{X}\mathbin{.}\overline{y-z}\mathbin{.}\overline{y-w}\mathbin{.}\overline{y-a}\mathbin{.}\overline{z-w}\mathbin{.}\overline{z-a}\mathbin{.}\overline{w-a}\\
&+ \sqrt{Y}\mathbin{.}\overline{z-w}\mathbin{.}\overline{z-a}\mathbin{.}\overline{z-x}\mathbin{.}\overline{w-a}\mathbin{.}\overline{w-x}\mathbin{.}\overline{a-x}\\
&+ \sqrt{Z}\mathbin{.}\overline{w-a}\mathbin{.}\overline{w-x}\mathbin{.}\overline{w-y}\mathbin{.}\overline{a-x}\mathbin{.}\overline{a-y}\mathbin{.}\overline{x-y}\\
&+ \sqrt{W}\mathbin{.}\overline{a-x}\mathbin{.}\overline{a-y}\mathbin{.}\overline{a-z}\mathbin{.}\overline{x-y}\mathbin{.}\overline{x-z}\mathbin{.}\overline{y-z}\bigr\},
\end{align*}
or as it may be written
\[
\Omega\mathbin{.}\overline{x-z}\mathbin{.}\overline{x-w}\mathbin{.}\overline{y-z}\mathbin{.}\overline{y-w} = \varepsilon\,\dfrac{\overline{a-z}\mathbin{.}\overline{a-w}}{x-y}\bigl[\overline{y-z}\mathbin{.}\overline{y-w}\mathbin{.}\overline{a-y}\,\sqrt{X} - \overline{x-z}\mathbin{.}\overline{x-w}\mathbin{.}\overline{a-x}\,\sqrt{Y}\bigr],
\]
\[
+\ \varepsilon\,\dfrac{\overline{a-x}\mathbin{.}\overline{a-y}}{z-w}\bigl[\overline{w-x}\mathbin{.}\overline{w-y}\mathbin{.}\overline{a-w}\,\sqrt{Z} - \overline{z-x}\mathbin{.}\overline{z-y}\mathbin{.}\overline{a-z}\,\sqrt{W}\bigr],
\]
an equation for the determination of $\Omega$.

Consider first the expression which multiplies $\varepsilon\mathbin{.}\overline{a-z}\mathbin{.}\overline{a-w}$; this is
\[
= \tfrac{1}{x-y}\bigl\{\overline{y-z}\mathbin{.}\overline{y-w}\mathbin{.}\overline{a-y}\,\sqrt{X} - \overline{x-z}\mathbin{.}\overline{x-w}\mathbin{.}\overline{a-x}\,\sqrt{Y}\bigr\}
\]
we have
\[
= -BE_{12}\, \tfrac{1}{x-y}\bigl\{\sqrt{b_1 e_1 f_1 a_2 c_2 d_2} - \sqrt{b_2 e_2 f_2 a_1 c_1 d_1}\bigr\},
\]
and multiplying this by $A_{12}\, B_{12}\, C_{12}\, D_{12}\, = \sqrt{a_1 c_1 d_1 a_2 c_2 d_2}$, we derive
\[
\sqrt{a^6 - \varepsilon^2}\, BE_{12}\mathbin{.}C_{12}\mathbin{.}D_{12}\mathbin{.}A_{12} = \tfrac{1}{x-y}\bigl\{c_2 d_2 a_1 \sqrt{X} - c_1 d_1 a_2 \sqrt{Y}\bigr\},
\]
and similarly two other equations; the system may be written
\begin{align*}
\sqrt{a^6-\varepsilon^2}\, BE\mathbin{.}C\mathbin{.}D\mathbin{.}A &= \tfrac{1}{x-y}\bigl\{c_2 d_2 a_1 \sqrt{X} - c_1 d_1 a_2 \sqrt{Y}\bigr\},\\
\sqrt{a^6-\varepsilon^2}\, CE\mathbin{.}D\mathbin{.}B\mathbin{.}A &= \tfrac{1}{x-y}\bigl\{d_2 b_2 a_1 \sqrt{X} - d_1 b_1 a_2 \sqrt{Y}\bigr\},\\
\sqrt{a^6-\varepsilon^2}\, DE\mathbin{.}B\mathbin{.}C\mathbin{.}A &= \tfrac{1}{x-y}\bigl\{b_2 c_2 a_1 \sqrt{X} - b_1 c_1 a_2 \sqrt{Y}\bigr\},
\end{align*}
the suffixes on the left-hand side being always $12$. The letters $b,c,d$ which enter cyclically into these equations are any three of the five letters other than $a$; the remaining two letters $e$ and $f$ enter symmetrically, for $BE$ is a mere abbreviation for the double triad $BEF\mathbin{.}ACD$; and the like for $CE$, and $DE$.

\bigskip

Multiplying these equations by
\[
\frac{b-z\mathbin{.}b-w}{b-c\mathbin{.}b-d},\quad \frac{c-z\mathbin{.}c-w}{c-d\mathbin{.}c-b},\quad \frac{d-z\mathbin{.}d-w}{d-b\mathbin{.}d-c},
\]
respectively, and then adding, the right-hand side becomes
\[
= \tfrac{1}{x-y}\bigl\{\overline{y-z}\mathbin{.}\overline{y-w}\mathbin{.}a_1\,\sqrt{X} - \overline{x-z}\mathbin{.}\overline{x-w}\mathbin{.}a_2\,\sqrt{Y}\bigr\}.
\]
Writing
\[
\frac{b-z\mathbin{.}b-w}{b-c\mathbin{.}b-d} = \frac{-1}{c-d\mathbin{.}d-b\mathbin{.}b-c}\cdot \overline{b-z}\mathbin{.}\overline{b-w}\mathbin{.}\overline{c-d},\ \text{etc.,}
\]
the left-hand side becomes
\[
\resizebox{\textwidth}{!}{$\displaystyle \frac{-\sqrt{a^6-\varepsilon^2}\, A}{c-d\mathbin{.}d-b\mathbin{.}b-c}\bigl\{\overline{c-d}\mathbin{.}\overline{b-z\mathbin{.}b-w}\mathbin{.}BE_{12}\mathbin{.}C_{12}\mathbin{.}D_{12} + \overline{d-b}\mathbin{.}\overline{c-z\mathbin{.}c-w}\mathbin{.}CE_{12}\mathbin{.}D_{12}\mathbin{.}B_{12} + \overline{b-c}\mathbin{.}\overline{d-z\mathbin{.}d-w}\mathbin{.}DE_{12}\mathbin{.}B_{12}\mathbin{.}C_{12}\bigr\}$},
\]
which for shortness may be written
\[
\frac{-\sqrt{a^6-\varepsilon^2}\, A}{c-d\mathbin{.}d-b\mathbin{.}b-c}\, \Sigma\bigl\{\overline{c-d}\mathbin{.}\overline{B_1'}\mathbin{.}BE_{12}\mathbin{.}C_{12}\mathbin{.}D_{12}\bigr\},
\]
the summation referring to the three terms obtained by the cyclical interchange of the letters $b,c,d$. The result thus is
\[
\tfrac{1}{x-y}\bigl\{\overline{y-z}\mathbin{.}\overline{y-w}\mathbin{.}a_1\sqrt{X} - \overline{x-z}\mathbin{.}\overline{x-w}\mathbin{.}a_2\sqrt{Y}\bigr\}
\]
\[
= \frac{-\sqrt{a^6-\varepsilon^2}\, A_{12}}{c-d\mathbin{.}d-b\mathbin{.}b-c}\,\Sigma\bigl\{\overline{c-d}\mathbin{.}\overline{B_1'}\mathbin{.}BE_{12}\mathbin{.}C_{12}\mathbin{.}D_{12}\bigr\}.
\]

Interchanging $x,y$ with $z,w$ respectively, we have of course to interchange the suffixes $1,2$ and $3,4$; we thus find
\[
\tfrac{1}{z-w}\bigl\{\overline{w-x}\mathbin{.}\overline{w-y}\mathbin{.}a_3 \sqrt{Z} - \overline{z-x}\mathbin{.}\overline{z-y}\mathbin{.}a_4 \sqrt{W}\bigr\}
\]
\[
= \frac{-\sqrt{a^6-\varepsilon^2}\, A_{34}}{c-d\mathbin{.}d-b\mathbin{.}b-c}\,\Sigma\bigl\{\overline{c-d}\mathbin{.}\overline{B_3'}\mathbin{.}BE_{34}\mathbin{.}C_{34}\mathbin{.}D_{34}\bigr\},
\]
and we hence find the value of $\Omega\mathbin{.}\overline{x-z}\mathbin{.}\overline{x-w}\mathbin{.}\overline{y-z}\mathbin{.}\overline{y-w}$. But $\Omega = \alpha a^3 + \beta a^2 + \gamma a + \delta$, is $= \sqrt{a^6 - \varepsilon^2}\mathbin{.}A_{12}\mathbin{.}A_{34}\mathbin{.}A_{56}$: the resulting equation divides by $A_{12}\mathbin{.}A_{34}$; throwing out this factor, we have
\[
\sqrt{a^6 - \varepsilon^2}\,(x-z\mathbin{.}x-w\mathbin{.}y-z\mathbin{.}y-w)(c-d\mathbin{.}d-b\mathbin{.}b-c)\, A_{56}
\]
\[
= A_{34}\,\Sigma\bigl\{\overline{c-d}\mathbin{.}\overline{B_1'}\mathbin{.}BE_{12}\mathbin{.}C_{12}\mathbin{.}D_{12}\bigr\} + A_{12}\,\Sigma\bigl\{\overline{c-d}\mathbin{.}\overline{B_3'}\mathbin{.}BE_{34}\mathbin{.}C_{34}\mathbin{.}D_{34}\bigr\},
\]
where, as before, the summations refer to the three terms obtained by the cyclical interchange of the letters $b,c,d$; these being any three of the five letters other than $a$; and the remaining two letters $e,f$ enter into the formula symmetrically. The formula gives thus for $A_{56}$ ten values which are of course equal to each other.

\bigskip

Writing for $a$ each letter in succession, we obtain formulae for each of the six single-letter functions $A_{56}$ of $p$ and $q$; and the factor
\[
\sqrt{a^6 - \varepsilon^2}\,(x-z\mathbin{.}x-w\mathbin{.}y-z\mathbin{.}y-w)
\]
is the same in all the formulae.

We require further the expressions for the double-letter functions of $p,q$. Considering for example the function $BE_{56}$, which is
\[
= \tfrac{1}{p-q}\bigl\{\sqrt{b_5 e_5 f_5 a_6 c_6 d_6} - \sqrt{b_6 e_6 f_6 a_5 c_5 d_5}\bigr\}.
\]
then multiplying by $DE_{56}\mathbin{.}A_{56}\mathbin{.}B_{56}\mathbin{.}C_{56}$ we have
\[
\sqrt{a^6 - \varepsilon^2}\, DE_{56}\mathbin{.}A_{56}\mathbin{.}B_{56}\mathbin{.}C_{56} = \tfrac{1}{p-q}\bigl\{a_5 b_5 c_5\,\sqrt{P} - a_6 b_6 c_6\,\sqrt{Q}\bigr\},
\]
\[
= \tfrac{1}{p-q}\bigl\{\overline{a-q}\mathbin{.}\overline{b-q}\mathbin{.}\overline{c-q}\,\sqrt{P} - \overline{a-p}\mathbin{.}\overline{b-p}\mathbin{.}\overline{c-p}\,\sqrt{Q}\bigr\},
\]
or recollecting that $\varepsilon\sqrt{P},\ \varepsilon\sqrt{Q}$ are $= \alpha p^3 + \beta p^2 + \gamma p + \delta$ and $\alpha q^3 + \beta q^2 + \gamma q + \delta$ respectively, this is
\[
\varepsilon\mathbin{.}DE_{56}\mathbin{.}A_{56}\mathbin{.}B_{56}\mathbin{.}C_{56}
\]
\[
= \tfrac{1}{p-q}\bigl\{\overline{a-q}\mathbin{.}\overline{b-q}\mathbin{.}\overline{c-q}\,(\alpha p^3 + \beta p^2 + \gamma p + \delta) - \overline{a-p}\mathbin{.}\overline{b-p}\mathbin{.}\overline{c-p}\,(\alpha q^3 + \beta q^2 + \gamma q + \delta)\bigr\}.
\]
Using the well-known identity
\begin{align*}
\alpha p^3 + \beta p^2 + \gamma p + \delta = &\ (\alpha a^3 + \beta a^2 + \gamma a + \delta)\frac{p-b\mathbin{.}p-c\mathbin{.}p-d}{b-a\mathbin{.}c-a\mathbin{.}d-a}\\
&+ (\alpha b^3 + \beta b^2 + \gamma b + \delta)\frac{p-c\mathbin{.}p-d\mathbin{.}p-a}{c-b\mathbin{.}d-b\mathbin{.}a-b}\\
&+ (\alpha c^3 + \beta c^2 + \gamma c + \delta)\frac{p-d\mathbin{.}p-a\mathbin{.}p-b}{d-c\mathbin{.}a-c\mathbin{.}b-c}\\
&+ (\alpha d^3 + \beta d^2 + \gamma d + \delta)\frac{p-a\mathbin{.}p-b\mathbin{.}p-c}{a-d\mathbin{.}b-d\mathbin{.}c-d}
\end{align*}
and the like expression for $\alpha q^3 + \beta q^2 + \gamma q + \delta$, there will be on the right-hand side terms involving
\[
\alpha a^3 + \beta a^2 + \gamma a + \delta,\quad \alpha b^3 + \beta b^2 + \gamma b + \delta,\quad \alpha c^3 + \beta c^2 + \gamma c + \delta,
\]
but the term in $\alpha d^3 + \beta d^2 + \gamma d + \delta$ will disappear of itself.

The term in $\alpha a^3 + \beta a^2 + \gamma a + \delta$ is
\[
\frac{1}{b-a\mathbin{.}c-a}\,(\alpha a^3 + \beta a^2 + \gamma a + \delta)\,\tfrac{1}{p-q}\bigl\{\overline{a-q}\mathbin{.}\overline{b-q}\mathbin{.}\overline{c-q}\,(\ )_p - \overline{a-p}\mathbin{.}\overline{b-p}\mathbin{.}\overline{c-p}\,(\ )_q\bigr\},
\]
where the expression in $(\ )$ is $= \overline{a-p}\mathbin{.}\overline{a-q}$: hence the term is
\[
\frac{\alpha a^3 + \beta a^2 + \gamma a + \delta}{b-a\mathbin{.}c-a}\,\overline{a-p}\mathbin{.}\overline{a-q}\mathbin{.}\overline{b-q}\mathbin{.}\overline{c-q}\mathbin{.}\overline{b-p}\mathbin{.}\overline{c-p},
\]
which is
\[
= \frac{\alpha a^3 + \beta a^2 + \gamma a + \delta}{b-a\mathbin{.}c-a}\,A_{56}^{\,2}\, B_{56}^{\,2}\, C_{56}^{\,2}\,.
\]
Forming the two other like terms, the equation is
\begin{align*}
\varepsilon\mathbin{.}DE_{56}\mathbin{.}A_{56}\mathbin{.}B_{56}\mathbin{.}C_{56} = &\ \frac{\alpha a^3 + \beta a^2 + \gamma a + \delta}{b-a\mathbin{.}c-a}\, A_{56}^{\,2} B_{56}^{\,2} C_{56}^{\,2}\\
&+ \frac{\alpha b^3 + \beta b^2 + \gamma b + \delta}{c-b\mathbin{.}a-b}\, A_{56}^{\,2} B_{56}^{\,2} C_{56}^{\,2}\\
&+ \frac{\alpha c^3 + \beta c^2 + \gamma c + \delta}{a-c\mathbin{.}b-c}\, A_{56}^{\,2} B_{56}^{\,2} C_{56}^{\,2}.
\end{align*}

But the expressions
\[
\alpha a^3 + \beta a^2 + \gamma a + \delta,\ \alpha b^3 + \beta b^2 + \gamma b + \delta,\ \alpha c^3 + \beta c^2 + \gamma c + \delta,
\]
are
\[
= \tfrac{\sqrt{a^6 - \varepsilon^2}}{\varepsilon}A_{12}\mathbin{.}A_{34}\mathbin{.}A_{56},\ \tfrac{\sqrt{a^6 - \varepsilon^2}}{\varepsilon}B_{12}\mathbin{.}B_{34}\mathbin{.}B_{56},\ \tfrac{\sqrt{a^6 - \varepsilon^2}}{\varepsilon}C_{12}\mathbin{.}C_{34}\mathbin{.}C_{56},
\]
respectively: the whole equation thus divides by $C_{56}\mathbin{.}A_{56}\mathbin{.}B_{56}$; throwing out this factor, and then multiplying each side by $-\dfrac{\sqrt{a^6 - \varepsilon^2}}{\varepsilon}$ we find
\[
\frac{\sqrt{a^6 - \varepsilon^2}}{\varepsilon}\, DE_{56} = \Sigma\frac{b-c\mathbin{.}A_{12}\mathbin{.}A_{34}\mathbin{.}B_{56}\mathbin{.}C_{56}}{b-c\mathbin{.}c-a\mathbin{.}a-b}\cdot\overline{(-\sqrt{a^6 - \varepsilon^2}/\varepsilon)},
\]
in which formula if we imagine
\[
-\,\dfrac{\sqrt{a^6 - \varepsilon^2}}{\varepsilon}\quad\text{and}\quad -\,\dfrac{\sqrt{a^6 - \varepsilon^2}}{\varepsilon}
\]
each replaced by its value in terms of the $xy$- and $zw$-functions, we have an equation of the form
\[
(x-z\mathbin{.}x-w\mathbin{.}y-z\mathbin{.}y-w)\, DE_{56} = \overline{x-z\mathbin{.}x-w\mathbin{.}y-z\mathbin{.}y-w}\cdot M,
\]
where $M$ is a given rational and integral function of the $16$ and $16$ functions $A_{12},AB_{12}$ and $A_{34},AB_{34}$ of $x$ and $y$ and of $z$ and $w$ respectively. The factor
\[
(x-z\mathbin{.}x-w\mathbin{.}y-z\mathbin{.}y-w)
\]
is retained on the left-hand side as being the same factor which enters into the equations for $A_{56}$, etc.: but on the right-hand side $x-z\mathbin{.}x-w\mathbin{.}y-z\mathbin{.}y-w$ should be expressed in terms of the $xy$- and $zw$-functions. This can be done by means of the identity
\[
x-z\mathbin{.}x-w\mathbin{.}y-z\mathbin{.}y-w = \Sigma\,\frac{\begin{vmatrix} 1 & x+y & xy \\ 1 & z+w & zw \\ 1 & a+b & ab \end{vmatrix}\begin{vmatrix} 1 & x+y & xy \\ 1 & z+w & zw \\ 1 & a+c & ac \end{vmatrix}}{a-b\mathbin{.}a-c},
\]
where the summation refers to the three terms obtained by the cyclical interchange of the letters $a,b,c$. The first determinant, multiplied by $a-b$, is in fact
\[
\begin{vmatrix} a-z\mathbin{.}a-w & a-x\mathbin{.}a-y \\ b-z\mathbin{.}b-w & b-x\mathbin{.}b-y \end{vmatrix},
\]
and the second determinant, multiplied by $-(a-c)$, is
\[
= \begin{vmatrix} a-z\mathbin{.}a-w & a-x\mathbin{.}a-y \\ c-z\mathbin{.}c-w & c-x\mathbin{.}c-y \end{vmatrix},
\]
so that the formula may also be written
\[
x-z\mathbin{.}x-w\mathbin{.}y-z\mathbin{.}y-w = \Sigma\,\frac{\begin{vmatrix} a-z\mathbin{.}a-w & a-x\mathbin{.}a-y \\ b-z\mathbin{.}b-w & b-x\mathbin{.}b-y \end{vmatrix}\begin{vmatrix} a-z\mathbin{.}a-w & a-x\mathbin{.}a-y \\ c-z\mathbin{.}c-w & c-x\mathbin{.}c-y \end{vmatrix}}{(a-b)^2(a-c)^2},
\]
or, what is the same thing, it is
\[
\overline{x-z\mathbin{.}x-w\mathbin{.}y-z\mathbin{.}y-w} = \Sigma\,\frac{\cdots}{\cdots},
\]
which is the required expression for $x-z\mathbin{.}x-w\mathbin{.}y-z\mathbin{.}y-w$; the letters $a,b,c$, which enter into the formula, are any three of the six letters.

As regards the verification of the identity, observe that it may be written
\[
x-z\mathbin{.}x-w\mathbin{.}y-z\mathbin{.}y-w = \Sigma\,\frac{\{L + M(a+b) + Nab\}\{L + M(a+c) + Nac\}}{a-b\mathbin{.}a-c},
\]
where $L,M,N$ are
\[
=(x+y)zw - (z+w)xy,\quad xy - zw,\quad\text{and}\quad z+w-x-y:
\]
this is readily reduced to
\[
x-z\mathbin{.}x-w\mathbin{.}y-z\mathbin{.}y-w = M^2 - NL,
\]
which can be at once verified.

\medskip
\hfill Cambridge, 12th March, 1879.

\bigskip

I take the opportunity of remarking that, in the double-letter formulae, the sign of the second term is, not as I have in general written it $-$, but is $+$,
\[
AB = \tfrac{1}{x-y}\bigl\{\sqrt{a_1 b_1 f_1 c_2 d_2 e_2} + \sqrt{a_2 b_2 f_2 c_1 d_1 e_1}\bigr\},\ \text{etc.}
\]
In fact, introducing a factor $\omega$ which is a function of $x$ and $y$, the odd and even $\vartheta$-functions are $= \omega\sqrt{a_1 a_2}$, etc., and
\[
= \omega\,\tfrac{1}{x-y}\bigl\{\sqrt{a_1 b_1 f_1 c_2 d_2 e_2} + \sqrt{a_2 b_2 f_2 c_1 d_1 e_1}\bigr\},\ \text{etc.,}
\]
respectively; $\omega$ is a function which on the interchange of $x,y$ changes only its sign; and this being so, then when $x$ and $y$ are interchanged, each single-letter function changes its sign, and each double-letter function remains unaltered.

\medskip
\hfill Cambridge, 29th July, 1879.

\bigskip\bigskip
\hrule
\bigskip

\begin{center}
\textbf{\large 704.}\\[4pt]
\textbf{A MEMOIR ON THE SINGLE AND DOUBLE THETA-FUNCTIONS.}
\end{center}

\medskip

\noindent\small [From the \textit{Philosophical Transactions of the Royal Society of London}, vol.\ 171, Part III., (1880), pp.\ 897--1002. Received November 14,---Read November 28, 1879.]\normalsize

\bigskip

The Theta-Functions, although arising historically from the Elliptic Functions, may be considered as in order of simplicity preceding these, and connecting themselves directly with the exponential function ($e^x$ or) $\exp.\ x$; viz.\ they may be defined each of them as a sum of a series of exponentials, singly infinite in the case of the single functions, doubly infinite in the case of the double functions; and so on. The number of the single functions is $4$; and the quotients of these, or say three of them each divided by the fourth, are the elliptic functions $\operatorname{sn},\operatorname{cn},\operatorname{dn}$; the number of the double functions is ($4^2 =$) $16$; and the quotients of these, or say fifteen of them each divided by the sixteenth, are the hyper-elliptic functions of two arguments depending on the square root of a sextic function. Generally, the number of the $p$-tuple theta-functions is $4^p$; and the quotients of these, or say all but one of them each divided by the remaining function, are the Abelian functions of $p$ arguments depending on the irrational function $y$ defined by the equation $F(x,y) = 0$ of a curve of deficiency $p$. If, instead of connecting the ratios of the functions with a plane curve, we consider the functions themselves as coordinates of a point in a space of $(4^p - 1)$ dimensions, then we have the single functions as the four coordinates of a point on a quadri-quadric curve (one-fold locus) in ordinary space; and the double functions as the sixteen coordinates of a point on a quadri-quadric two-fold locus in $15$-dimensional space, the deficiency of this two-fold locus being of course $= 2$.

The investigations contained in the First Part of the present Memoir, although for simplicity of notation exhibited only in regard to the double functions are, in fact, applicable to the general case of the $p$-tuple functions; but in the main the Memoir relates only to the single and double functions, and the title has been given to it accordingly. The investigations just referred to extend to the single functions and there is, it seems to me, an advantage in carrying on the two theories simultaneously up to and inclusive of the establishment of what I call the Product-theorem: this is a natural point of separation for the theories of the single and the double functions respectively. The ulterior developments of the two theories are indeed closely analogous to each other; but on the one hand the course of the single theory would be only with difficulty perceptible in the greater complexity of the double theory; and on the other hand we need the single theory as a guide for the course of the double theory.

I accordingly stop to point out in a general manner the course of the single theory, and, in connexion with it but more briefly, that of the double theory; and I then, in the Second and the Third Parts respectively, consider in detail the two theories separately; first, that of the single functions, and then that of the double functions. The paragraphs of the Memoir are numbered consecutively.

The definition adopted for the theta-functions differs somewhat from that which is ordinarily used.

The earlier memoirs on the double theta-functions are the well-known ones:

\medskip

Rosenhain, ``M\'emoire sur les fonctions de deux variables et \`a quatre p\'eriodes, qui sont les inverses des int\'egrales ultra-elliptiques de la premi\`ere classe.'' [1846.] Paris: --- M\'em.\ Savans \'Etrang., t.\ XI.\ (1851), pp.\ 361--468.

\medskip

G\"opel, ``Theoriae transcendentium Abelianarum primi ordinis adumbratio levis,'' Crelle, t.\ XXXV.\ (1847), pp.\ 277--312.

\medskip

My first paper --- Cayley, ``On the Double $\vartheta$-Functions in connexion with a 16-nodal Surface,'' Crelle-Borchardt, t.\ LXXXIII.\ (1877), pp.\ 210--219, [662]---was founded directly upon these, and was immediately followed by Dr Borchardt's paper,

\medskip

Borchardt, ``Ueber die Darstellung der Kummer'sche Fl\"ache vierter Ordnung mit sechzehn Knotenpunkten durch die G\"opelsche biquadratische Relation zwischen vier Thetafunctionen mit zwei Variabeln,'' Ditto, pp.\ 234--244.

\medskip

My other later papers, [663, 664, 665, 697, 703], are contained in the same Journal.

\bigskip

\begin{center}\textbf{FIRST PART.---INTRODUCTORY.}\end{center}

\medskip

\textit{Definition of the theta-functions.}

\medskip

\textbf{1.} The $p$-tuple functions depend upon $\tfrac{1}{2}p(p+1)$ parameters which are the coefficients of a quadric function of $p$ ultimately disappearing integers, upon $p$ arguments, and upon $2p$ characters, each $=0$ or $1$, which form the characteristic of the $4^p$ functions; but it will be sufficient to write down the formulae in the case $p = 2$.

As already mentioned, the adopted definition differs somewhat from that which is ordinarily used. I use, as will be seen, a quadric function $\tfrac{1}{4}(a,h,b\mid m,n)^2$ with even integer values of $m,n$, instead of $(a,h,b\mid m,n)^2$ with even or odd values; and I write the other term $\tfrac{1}{2}\pi i\,(mu + nv)$, instead of $mu + nv$; this comes to affecting the arguments $u,v$ with a factor $\tfrac{1}{2}\pi i$, so that the quarter-periods (instead of being $\tfrac{1}{2}\pi i$) are made to be $= 1$.

\medskip

\textbf{2.} We write
\[
\bigl(\ \bigr) = \tfrac{1}{4}(a,h,b\mid m,n)^2 + \tfrac{1}{2}\pi i\,(mu + nv),
\]
and in like manner
\[
\Bigl(\!\!\bigm|\!\!^{m+\alpha}_{n+\beta}\Bigm|\!\!^{u+\gamma}_{v+\delta}\Bigr) = \tfrac{1}{4}(a,h,b \mid m + \alpha, n + \beta)^2 + \tfrac{1}{2}\pi i\bigl\{(m + \alpha)(u + \gamma) + (n + \beta)(v + \delta)\bigr\},
\]
and prefixing to either of these the functional symbol $\exp.$ we have the exponential of the function in question, that is, $e$ with the function as an exponent.

We then write, as the definition of the double theta-functions,
\[
\vartheta\!\begin{pmatrix} \alpha,\beta \\ \gamma,\delta \end{pmatrix}\!(u,v) = \Sigma\exp.\Bigl(\!\!\bigm|\!\!^{m+\alpha}_{n+\beta}\Bigm|\!\!^{u+\gamma}_{v+\delta}\Bigr)
\]
where the summation extends to all positive and negative even integer values (zero included) of $m$ and $n$ respectively; $\alpha,\beta,\gamma,\delta$ might denote any quantities whatever, but for the theta-functions they are regarded as denoting positive or negative integers; this being so, it will appear that the only effect of altering each or any of them by an even integer is to reverse (it may be) the sign of the function; and the distinct functions are consequently the ($4^2 =$) $16$ functions obtained by giving to each of the quantities $\alpha,\beta,\gamma,\delta$ the two values $0$ and $1$ successively.

\medskip

\textbf{3.} We thus have the double theta-functions, depending on the parameters $(a,h,b)$ which determine the quadric function $(a,h,b\mid m,n)^2$ of the disappearing even integers $(m,n)$, and on the two arguments $(u,v)$: in the symbol $\begin{pmatrix} \alpha,\beta \\ \gamma,\delta \end{pmatrix}$, which is called the characteristic, the characters $\alpha,\beta,\gamma,\delta$ are each of them $= 0$ or $1$; and we thus have the $16$ functions.

\medskip

\textbf{4.} The parameters $(a,h,b)$ may be real or imaginary, but they must be such that reducing each of them to its real part the resulting function $(\ast\mid m,n)^2$ is invariable in its sign, and negative for all real values of $m$ and $n$: this is, in fact, the condition for the convergency of the series which give the values of the theta-functions.

The characteristic $\begin{pmatrix} \alpha,\beta \\ \gamma,\delta \end{pmatrix}$ is said to be even or odd according as the sum $\alpha\gamma + \beta\delta$ is even or odd.

\medskip

\textit{Allied functions.}

\medskip

\textbf{5.} As already remarked, the definition of
\[
\vartheta\!\begin{pmatrix} \alpha,\beta \\ \gamma,\delta \end{pmatrix}\!(u,v)
\]
is not restricted to the case where the $\alpha,\beta,\gamma,\delta$ represent integers, and there is actually occasion to consider functions of this form where they are not integers: in particular, $\alpha,\beta$ may be either or each of them of the form, integer $+ \tfrac{1}{2}$. But the functions thus obtained are not regarded as \textit{theta-functions}, and the expression theta-function will consequently not extend to include them.

\medskip

\textit{Properties of the Theta-Functions: Various sub-headings.}

\medskip

\textit{Even-integer alteration of characters.}

\medskip

\textbf{6.} If $x,y$ be integers, then $m,n$ having the several even integer values from $-\infty$ to $+\infty$ respectively, it is obvious that $m + \alpha + 2x,\ n + \beta + 2y$ will have the same series of values with $m + \alpha,\ n + \beta$ respectively; and it thence follows that
\[
\vartheta\!\begin{pmatrix} \alpha + 2x,\beta + 2y \\ \gamma,\delta \end{pmatrix}\!(u,v) = \vartheta\!\begin{pmatrix} \alpha,\beta \\ \gamma,\delta \end{pmatrix}\!(u,v).
\]
Similarly if $z,w$ are integers, then in the function
\[
\vartheta\!\begin{pmatrix} \alpha,\beta \\ \gamma + 2z,\delta + 2w \end{pmatrix}\!(u,v)
\]
the argument of the exponential function contains the term
\[
\tfrac{1}{2}\pi i\bigl\{m + \alpha\mathbin{.}u + \gamma + 2z + n + \beta\mathbin{.}v + \delta + 2w\bigr\}
\]
this differs from its original value by
\[
\tfrac{1}{2}\pi i\,(m + \alpha\mathbin{.}2z + n + \beta\mathbin{.}2w),
\]
\[
= \pi i\,(mz + nw) + \pi i\,(\alpha z + \beta w),
\]
and then, $m$ and $n$ being even integers, $mz + nw$ is also an even integer, and the term $\pi i\,(mz + nw)$ does not affect the value of the exponential: we thus introduce into each term of the series the factor $\exp.\ \pi i\,(\alpha z + \beta w)$, which is, in fact, $(-)^{\alpha z + \beta w}$, and we consequently have
\[
\vartheta\!\begin{pmatrix} \alpha,\beta \\ \gamma + 2z,\delta + 2w \end{pmatrix}\!(u,v) = (-)^{\alpha z + \beta w}\,\vartheta\!\begin{pmatrix} \alpha,\beta \\ \gamma,\delta \end{pmatrix}\!(u,v),
\]
or, uniting the two results,
\[
\vartheta\!\begin{pmatrix} \alpha + 2x,\beta + 2y \\ \gamma + 2z,\delta + 2w \end{pmatrix}\!(u,v) = (-)^{\alpha z}\,(-)^{\beta w}\,\vartheta\!\begin{pmatrix} \alpha,\beta \\ \gamma,\delta \end{pmatrix}\!(u,v).
\]
This sustains the before-mentioned conclusion that the only distinct functions are the $16$ functions obtained by giving to the characters $\alpha,\beta,\gamma,\delta$ the values $0$ and $1$ respectively.

\medskip

\textit{Odd-integer alteration of characters.}

\medskip

\textbf{7.} The effect is obviously to interchange the different functions.

\medskip

\textit{Even and odd functions.}

\medskip

\textbf{8.} It is clear that $-m - \alpha,\ -n - \beta$ have precisely the same series of values with $m + \alpha,\ n + \beta$ respectively: hence considering the function
\[
\vartheta\!\begin{pmatrix} \alpha,\beta \\ \gamma,\delta \end{pmatrix}\!(-u, -v),
\]
the linear term in the argument of the exponential may be taken to be
\[
\tfrac{1}{2}\pi i\bigl\{-m - \alpha\mathbin{.}-u + \gamma + -n - \beta\mathbin{.}-v + \delta\bigr\},
\]
which is
\[
= \tfrac{1}{2}\pi i\bigl\{m + \alpha\mathbin{.}u + \gamma + n + \beta\mathbin{.}v + \delta\bigr\} - \pi i\bigl\{m + \alpha\mathbin{.}\gamma + n + \beta\mathbin{.}\delta\bigr\};
\]
the second term is here
\[
= -\pi i\,(m\gamma + n\delta) - \pi i\,(\alpha\gamma + \beta\delta),
\]
where, $m\gamma + n\delta$ being an even integer, the part $-\pi i\,(m\gamma + n\delta)$ does not alter the value of the exponential: the effect of the remaining part $-\pi i\,(\alpha\gamma + \beta\delta)$ is to affect each term of the series with the factor $\exp.\ -\pi i\,(\alpha\gamma + \beta\delta)$, or what is the same thing, $\exp.\ \pi i\,(\alpha\gamma + \beta\delta)$, each of these being, in fact, $(-)^{\alpha\gamma + \beta\delta}$.

\medskip

Viz.\ $\vartheta\!\begin{pmatrix} \alpha,\beta \\ \gamma,\delta \end{pmatrix}\!(u,v)$ is an even or odd function of the two arguments $(u,v)$ conjointly, according as the characteristic $\begin{pmatrix} \alpha,\beta \\ \gamma,\delta \end{pmatrix}$ is even or odd.

\medskip

\textit{The quarter-periods unity.}

\medskip

\textbf{9.} Taking $z$ and $w$ integers, we have from the definition
\[
\vartheta\!\begin{pmatrix} \alpha,\beta \\ \gamma,\delta \end{pmatrix}\!(u + z,\, v + w) = \vartheta\!\begin{pmatrix} \alpha,\beta \\ \gamma + z,\delta + w \end{pmatrix}\!(u,v),
\]
viz.\ the effect of altering the arguments $u,v$ into $u + z,\ v + w$ is simply to interchange the functions as shown by this formula.

If $z$ and $w$ are each of them even, then replacing them by $2z, 2w$ respectively, we have
\[
\vartheta\!\begin{pmatrix} \alpha,\beta \\ \gamma,\delta \end{pmatrix}\!(u + 2z,\, v + 2w) = \vartheta\!\begin{pmatrix} \alpha,\beta \\ \gamma + 2z,\delta + 2w \end{pmatrix}\!(u,v),
\]
which by a preceding formula is
\[
= (-)^{\alpha z + \beta w}\,\vartheta\!\begin{pmatrix} \alpha,\beta \\ \gamma,\delta \end{pmatrix}\!(u,v),
\]
or the function is altered at most in its sign. And again writing $2z, 2w$ for $z,w$, we have
\[
\vartheta\!\begin{pmatrix} \alpha,\beta \\ \gamma,\delta \end{pmatrix}\!(u + 4z,\, v + 4w) = \vartheta\!\begin{pmatrix} \alpha,\beta \\ \gamma,\delta \end{pmatrix}\!(u,v).
\]
In reference to the foregoing results we say that the theta-functions have the quarter-periods $(1,1)$, the half-periods $(2,2)$, and the whole periods $(4,4)$.

\medskip

\textit{The conjoint quarter quasi-periods.}

\medskip

\textbf{10.} Taking $x,y$ integers, we consider the effect of the change of $u,v$ into
\[
u + \tfrac{1}{\pi i}(ax + hy),\quad v + \tfrac{1}{\pi i}(hx + by).
\]
It is convenient to start from the function
\[
\vartheta\!\begin{pmatrix} \alpha - x,\beta - y \\ \gamma,\delta \end{pmatrix}\!\bigl(u + \tfrac{1}{\pi i}(ax + hy),\, v + \tfrac{1}{\pi i}(hx + by)\bigr);
\]
the argument of the exponential is here
\[
\tfrac{1}{4}(a,h,b\mid m + \alpha - x, n + \beta - y)^2
\]
\[
+ \tfrac{1}{2}\pi i\bigl\{m + \alpha - x\mathbin{.}u + \gamma + \tfrac{1}{\pi i}(ax + hy) + n + \beta - y\mathbin{.}v + \delta + \tfrac{1}{\pi i}(hx + by)\bigr\},
\]
which is
\[
= \tfrac{1}{4}(a,h,b\mid m + \alpha, n + \beta)^2 + \tfrac{1}{2}\pi i\,(m + \alpha\mathbin{.}u + \gamma + n + \beta\mathbin{.}v + \delta)
\]
$+$ other terms which are as follows: viz.\ they are
\[
-\tfrac{1}{2}(a,h,b\mid m + \alpha, n + \beta\mid x,y) + \tfrac{1}{2}(m + \alpha\mathbin{.}ax + hy + n + \beta\mathbin{.}hx + by)
\]
\[
+ \tfrac{1}{4}(a,h,b\mid x,y)^2 - \tfrac{1}{2}\pi i\,(x\mathbin{.}u + \gamma + y\mathbin{.}v + \delta)
\]
\[
- \tfrac{1}{2}(x\mathbin{.}ax + hy + y\mathbin{.}hx + by),
\]
where the terms of the right-hand column are, in fact,
\[
= \tfrac{1}{2}\pi i\,(m + \alpha\mathbin{.}ax + hy + n + \beta\mathbin{.}hx + by) - \tfrac{1}{2}(a,h,b\mid m + \alpha, n + \beta\mid x,y)
\]
\[
-\tfrac{1}{2}\pi i\,(x\mathbin{.}u + \gamma + y\mathbin{.}v + \delta) - \tfrac{1}{4}(a,h,b\mid x,y)^2,
\]
and the other terms in question thus reduce themselves to
\[
-\tfrac{1}{4}(a,h,b\mid x,y)^2 - \tfrac{1}{2}\pi i\,(x\mathbin{.}u + \gamma + y\mathbin{.}v + \delta),
\]
which are independent of $m,n$, and they thus affect each term of the series with the same exponential factor. The result is
\[
\vartheta\!\begin{pmatrix} \alpha - x,\beta - y \\ \gamma,\delta \end{pmatrix}\!\bigl(u + \tfrac{1}{\pi i}(ax + hy),\, v + \tfrac{1}{\pi i}(hx + by)\bigr)
\]
\[
= \exp.\bigl\{-\tfrac{1}{4}(a,h,b\mid x,y)^2 - \tfrac{1}{2}\pi i\,(x\mathbin{.}u + \gamma + y\mathbin{.}v + \delta)\bigr\}\,\vartheta\!\begin{pmatrix} \alpha,\beta \\ \gamma,\delta \end{pmatrix}\!(u,v);
\]
or (what is the same thing) for $\alpha,\beta$ writing $\alpha + x, \beta + y$ respectively, we have
\[
\vartheta\!\begin{pmatrix} \alpha,\beta \\ \gamma,\delta \end{pmatrix}\!\bigl(u + \tfrac{1}{\pi i}(ax + hy),\, v + \tfrac{1}{\pi i}(hx + by)\bigr)
\]
\[
= \exp.\bigl\{-\tfrac{1}{4}(a,h,b\mid x,y)^2 - \tfrac{1}{2}\pi i\,(x\mathbin{.}u + \gamma + y\mathbin{.}v + \delta)\bigr\}\,\vartheta\!\begin{pmatrix} \alpha + x,\beta + y \\ \gamma,\delta \end{pmatrix}\!(u,v).
\]
Taking $x,y$ even, or writing $2x, 2y$ for $x,y$, then on the right-hand side we have
\[
\vartheta\!\begin{pmatrix} \alpha + 2x,\beta + 2y \\ \gamma,\delta \end{pmatrix}\!(u,v),
\]
which is
\[
= \vartheta\!\begin{pmatrix} \alpha,\beta \\ \gamma,\delta \end{pmatrix}\!(u,v):
\]
but there is still the exponential factor.

\medskip

\textbf{11.} The formulae show that the effect of the change $u,v$ into $u + \tfrac{1}{\pi i}(ax + hy)$, $v + \tfrac{1}{\pi i}(hx + by)$, where $x,y$ are integers, is to interchange the functions, affecting them however with an exponential factor; and we hence say that $\tfrac{1}{\pi i}(a,h)$, $\tfrac{1}{\pi i}(h,b)$ are conjoint quarter quasi-periods.

\medskip

\textit{The product-theorem.}

\medskip

\textbf{12.} We multiply two theta-functions
\[
\vartheta\!\begin{pmatrix} \alpha,\beta \\ \gamma,\delta \end{pmatrix}\!(u,v),\quad \vartheta\!\begin{pmatrix} \alpha',\beta' \\ \gamma',\delta' \end{pmatrix}\!(u',v');
\]
it is found that the result is a sum of four products
\[
\Theta\!\begin{pmatrix} \tfrac{1}{2}(\alpha + \alpha') + p,\tfrac{1}{2}(\beta + \beta') + q \\ \gamma + \gamma',\delta + \delta' \end{pmatrix}\!(2u, 2v)\cdot \Theta\!\begin{pmatrix} \tfrac{1}{2}(\alpha - \alpha') + p,\tfrac{1}{2}(\beta - \beta') + q \\ \gamma - \gamma',\delta - \delta' \end{pmatrix}\!(2u', 2v'),
\]
where $p,q$ have in the four products respectively the values $(0,0),(1,0),(0,1)$, and $(1,1)$; $\Theta$ is written in place of $\vartheta$ to denote that the parameters $(a,h,b)$ are to be changed into $(2a, 2h, 2b)$. It is to be noticed that, if $\alpha,\alpha'$ are both even or both odd, then $\tfrac{1}{2}(\alpha + \alpha'), \tfrac{1}{2}(\alpha - \alpha')$ are integers; and so, if $\beta,\beta'$ are both even or both odd, then $\tfrac{1}{2}(\beta + \beta'), \tfrac{1}{2}(\beta - \beta')$ are integers; and these conditions being satisfied (and in particular they are so if $\alpha = \alpha',\beta = \beta'$) then the functions on the right-hand side of the equation are theta-functions (with new parameters as already mentioned); but if the conditions are not satisfied, then the functions on the right-hand side are only allied functions. In the applications of the theorem the functions on the right-hand side are eliminated between the different equations, as will appear.

\medskip

\textbf{13.} The proof is immediate: in the first of the theta-functions, the argument of the exponential is
\[
\Bigl(\!\!\bigm|\!\!^{m+\alpha}_{n+\beta}\Bigm|\!\!^{u+\gamma}_{v+\delta}\Bigr),
\]
and in the second, writing $m',n'$ instead of $m,n$, the argument is
\[
\Bigl(\!\!\bigm|\!\!^{m'+\alpha'}_{n'+\beta'}\Bigm|\!\!^{u'+\gamma'}_{v'+\delta'}\Bigr);
\]
hence in the product, the argument of the exponential is the sum of these two functions, viz.
\[
= \tfrac{1}{4}(a,h,b\mid m + \alpha, n + \beta)^2 + \tfrac{1}{2}\pi i\,(m + \alpha\mathbin{.}u + \gamma + n + \beta\mathbin{.}v + \delta)
\]
\[
+ \tfrac{1}{4}(a,h,b\mid m' + \alpha', n' + \beta')^2 + \tfrac{1}{2}\pi i\,(m' + \alpha'\mathbin{.}u' + \gamma' + n' + \beta'\mathbin{.}v' + \delta').
\]
Comparing herewith the sum of the two functions
\[
\Bigl(\!\!\bigm|\!\!^{\mu + \tfrac{1}{2}(\alpha + \alpha')}_{\nu + \tfrac{1}{2}(\beta + \beta')}\Bigm|\!\!^{2u + \gamma + \gamma'}_{2v + \delta + \delta'}\Bigr)',\quad \Bigl(\!\!\bigm|\!\!^{\mu' + \tfrac{1}{2}(\alpha - \alpha')}_{\nu' + \tfrac{1}{2}(\beta - \beta')}\Bigm|\!\!^{2u' + \gamma - \gamma'}_{2v' + \delta - \delta'}\Bigr)',
\]
\[
= \tfrac{1}{4}(2a, 2h, 2b\mid \mu + \tfrac{1}{2}(\alpha + \alpha'), \nu + \tfrac{1}{2}(\beta + \beta'))^2
\]
\[
\hspace{2em}+ \tfrac{1}{2}\pi i\bigl\{\mu + \tfrac{1}{2}(\alpha + \alpha')\mathbin{.}2u + \gamma + \gamma' + \nu + \tfrac{1}{2}(\beta + \beta')\mathbin{.}2v + \delta + \delta'\bigr\}
\]
\[
+ \tfrac{1}{4}(2a, 2h, 2b\mid \mu' + \tfrac{1}{2}(\alpha - \alpha'), \nu' + \tfrac{1}{2}(\beta - \beta'))^2
\]
\[
\hspace{2em}+ \tfrac{1}{2}\pi i\bigl\{\mu' + \tfrac{1}{2}(\alpha - \alpha')\mathbin{.}2u' + \gamma - \gamma' + \nu' + \tfrac{1}{2}(\beta - \beta')\mathbin{.}2v' + \delta - \delta'\bigr\}
\]
the two sums are identical if only
\[
m + m' = 2\mu,\quad n + n' = 2\nu,\quad m - m' = 2\mu',\quad n - n' = 2\nu',
\]
as may easily be verified by comparing the quadric and the linear terms separately. The product of the two theta-functions is thus
\[
\Sigma\exp.\Bigl(\!\!\bigm|\!\!^{\mu + \tfrac{1}{2}(\alpha + \alpha')}_{\nu + \tfrac{1}{2}(\beta + \beta')}\Bigm|\!\!^{2u + \gamma + \gamma'}_{2v + \delta + \delta'}\Bigr)'\cdot \Sigma\exp.\Bigl(\!\!\bigm|\!\!^{\mu' + \tfrac{1}{2}(\alpha - \alpha')}_{\nu' + \tfrac{1}{2}(\beta - \beta')}\Bigm|\!\!^{2u' + \gamma - \gamma'}_{2v' + \delta - \delta'}\Bigr)',
\]
with the proper conditions as to the values of $\mu,\nu$ and of $\mu',\nu'$ in the two sums respectively. As to this, observe that $m,m'$ are even integers; say for a moment that they are \textit{similar} when they are both $\equiv 0$ or both $\equiv 2\ (\modd 4)$, but \textit{dissimilar} when they are one of them $\equiv 0$ and the other of them $\equiv 2\ (\modd 4)$; and the like as regards $n,n'$. Hence if $m,m'$ are similar, $\mu,\mu'$ are both of them even; but if $m,m'$ are dissimilar, then $\mu,\mu'$ are both of them odd. And so if $n,n'$ are similar, $\nu,\nu'$ are both of them even; but if $n,n'$ are dissimilar, then $\nu,\nu'$ are both odd.

\medskip

\textbf{14.} There are four cases:
\begin{align*}
&m,m'\ \text{similar},\quad n,n'\ \text{similar},\\
&m,m'\ \text{dissimilar},\quad n,n'\ \text{similar},\\
&m,m'\ \text{similar},\quad n,n'\ \text{dissimilar},\\
&m,m'\ \text{dissimilar},\quad n,n'\ \text{dissimilar}.
\end{align*}
In the first of these, $\mu,\nu,\mu',\nu'$ are all of them even, and the product is
\[
= \Theta\!\begin{pmatrix} \tfrac{1}{2}(\alpha + \alpha'),\tfrac{1}{2}(\beta + \beta') \\ \gamma + \gamma',\delta + \delta' \end{pmatrix}\!(2u, 2v)\cdot \Theta\!\begin{pmatrix} \tfrac{1}{2}(\alpha - \alpha'),\tfrac{1}{2}(\beta - \beta') \\ \gamma - \gamma',\delta - \delta' \end{pmatrix}\!(2u', 2v').
\]
In the second case, writing $\mu + 1,\mu' + 1$ for $\mu,\mu'$, the new values of $\mu,\mu'$ will be both even, and we have the like expression with only the characters $\tfrac{1}{2}(\alpha + \alpha'),\tfrac{1}{2}(\alpha - \alpha')$ each increased by $1$; so in the third case we obtain the like expression with only the characters $\tfrac{1}{2}(\beta + \beta'),\tfrac{1}{2}(\beta - \beta')$ each increased by $1$; and in the fourth case the like expression with the four upper characters each increased by $1$. The product of the two theta-functions is thus equal to the sum of the four products, according to the theorem.

\medskip

\textit{R\'esum\'e of the ulterior theory of the single functions.}

\medskip

\textbf{15.} For the single theta-functions the Product-theorem comprises $16$ equations, and for the double theta-functions, $256$ equations: these systems will be given in full in the sequel. But attending at present to the single functions, I write down here the first four of the $16$ equations, viz.\ these are
\begin{align*}
0\mathbin{.}0\quad &\vartheta\!\begin{pmatrix} 0 \\ 0 \end{pmatrix}\!(u + u')\,\vartheta\!\begin{pmatrix} 0 \\ 0 \end{pmatrix}\!(u - u') = XX' + YY',\\
1\mathbin{.}0\quad &\vartheta\!\begin{pmatrix} 1 \\ 0 \end{pmatrix}\,,,\quad \vartheta\!\begin{pmatrix} 1 \\ 0 \end{pmatrix},, = YX' + XY',\\
0\mathbin{.}1\quad &\vartheta\!\begin{pmatrix} 0 \\ 1 \end{pmatrix}\,,,\quad \vartheta\!\begin{pmatrix} 0 \\ 1 \end{pmatrix},, = XX' - YY',\\
1\mathbin{.}1\quad &\vartheta\!\begin{pmatrix} 1 \\ 1 \end{pmatrix}\,,,\quad \vartheta\!\begin{pmatrix} 1 \\ 1 \end{pmatrix},, = -YX' + XY';
\end{align*}
where $X,Y$ denote $\Theta\!\begin{pmatrix} 0 \\ 0 \end{pmatrix}\!(2u), \Theta\!\begin{pmatrix} 1 \\ 0 \end{pmatrix}\!(2u)$ respectively, and $X',Y'$ the same functions of $2u'$ respectively. In the other equations we have on the left-hand the product of different theta-functions of $u + u',\ u - u'$ respectively, and on the right-hand expressions involving other functions, $X_1, Y_1, X_1', Y_1'$, \&c., of $2u$ and $2u'$ respectively.

\medskip

\textbf{16.} By writing $u' = 0$, we have on the left-hand, squares or products of theta-functions of $u$, and on the right-hand expressions containing functions of $2u$: in particular, the above equations show that the squares of the four theta-functions are equal to linear functions of $X,Y$; that is, there exist between the squared functions two linear relations: or again, introducing a variable argument $x$, the four squared functions may be taken to be proportional to linear functions
\[
\mathfrak{A}(a - x),\ \mathfrak{B}(b - x),\ \mathfrak{C}(c - x),\ \mathfrak{D}(d - x),
\]
where $\mathfrak{A},\mathfrak{B},\mathfrak{C},\mathfrak{D},\,a,b,c,d$ are constants. This suggests a new notation for the four functions, viz.\ we write
\[
\vartheta\!\begin{pmatrix} 0 \\ 0 \end{pmatrix}\!(u) = Au,\quad \vartheta\!\begin{pmatrix} 1 \\ 0 \end{pmatrix}\!(u) = Bu,\quad \vartheta\!\begin{pmatrix} 0 \\ 1 \end{pmatrix}\!(u) = Cu,\quad \vartheta\!\begin{pmatrix} 1 \\ 1 \end{pmatrix}\!(u) = Du;
\]
and the result just mentioned then is
\[
A^2 u : B^2 u : C^2 u : D^2 u = \mathfrak{A}(a - x) : \mathfrak{B}(b - x) : \mathfrak{C}(c - x) : \mathfrak{D}(d - x),
\]
which expresses that the four functions are the coordinates of a point on a quadri-quadric curve in ordinary space.

\medskip

\textbf{17.} The remaining $12$ of the $16$ equations then contain on the left-hand products such as
\[
A(u + u')\mathbin{.}B(u - u');
\]
and by suitably combining them we obtain equations such as
\[
\frac{B(u + u')A(u - u') - A(u + u')B(u - u')}{D(u + u')C(u - u') + C(u + u')D(u - u')} = \text{function }(u'),
\]
where for brevity the arguments are written above; viz.\ the numerator of the fraction is
\[
B(u + u')\,A(u - u') - A(u + u')\,B(u - u'),
\]
and its denominator is
\[
D(u + u')\,C(u - u') + C(u + u')\,D(u - u').
\]
Admitting the form of the equation, the value of the function of $u'$ is at once found by writing in the equation $u = 0$; it is, as it ought to be, a function vanishing for $u' = 0$.

\medskip

\textbf{18.} Take in this equation $u$ indefinitely small; each side divides by $u'$, and the resulting equation is
\[
\frac{A'u\, Bu - Bu\, A'u}{C_0 u\, D_0 u} = \text{const.},
\]
where $A'u, B'u$ are the derived functions, or differential coefficients in regard to $u$. It thus appears that the combination $A'u\, Bu - Bu\, A'u$ is a constant multiple of $Cu\, Du$: or, what is the same thing, that the differential coefficient of the quotient-function $\dfrac{B}{A}u$ is a constant multiple of the product of the two quotient-functions $\dfrac{C}{A}u$ and $\dfrac{D}{A}u$.

\medskip

\textbf{19.} And then substituting for the several quotient-functions their values in terms of $x$, we obtain a differential relation between $x,u$; viz.\ the form hereof is
\[
\frac{du}{1} = \frac{M\,dx}{\sqrt{a-x\mathbin{.}b-x\mathbin{.}c-x\mathbin{.}d-x}}
\]
and it thus appears that the quotient-functions are in fact elliptic-functions: the actual values as obtained in the sequel are
\begin{align*}
\operatorname{sn} Ku &= \frac{Du}{Bu} = \sqrt{\frac{k'}{k}}\,\frac{Du}{Cu},\\
\operatorname{cn} Ku &= \sqrt{\frac{k'}{k}}\,\frac{Bu}{Cu},\\
\operatorname{dn} Ku &= \sqrt{k'}\,\frac{Au}{Cu};
\end{align*}
and we thus of course identify the functions $Au, Bu, Cu, Du$ with the $H$ and the $\Theta$ functions of Jacobi.

\medskip

\textbf{20.} If in the above-mentioned four equations we write first $u' = 0$, and then $= u' = 0$, and by means of the results eliminate from the original equations the quantities $X,Y,X',Y'$ which occur therein, we obtain expressions for the four products such as $A(u + u')\,A(u - u')$. One of these equations is
\[
C\!_0\, 0\mathbin{.}C(u + u')\,C(u - u') = C^2 u\, C^2 u' - D^2 u\, D^2 u'.
\]
Taking herein $u'$ indefinitely small, we obtain
\[
\frac{Cu\,C''u - (C'u)^2}{Cu} = \frac{C''_0}{C_0} - \left(\frac{D_0}{C_0}\right)^2 \frac{D^2 u}{C^2 u},
\]
where the left-hand side is in fact $\dfrac{d^2}{du^2}\log Cu$, or this second derived function of the theta-function $Cu$ is given in terms of the quotient-function $\dfrac{D}{C}u$; hence, integrating twice and taking the exponential of each side, we obtain $Cu$ as an exponential the argument of which contains the double integral $\int\!\int\!\left(\dfrac{Du}{Cu}\right)^2(du)^2$, of a squared quotient-function. This, in fact, corresponds to Jacobi's equation.

\medskip

\textbf{21.} From the same equation
\[
C\!_0\,0\mathbin{.}C(u + u')\,C(u - u') = C^2 u\, C^2 u' - D^2 u\, D^2 u',
\]
differentiating logarithmically in regard to $u'$ and integrating in regard to $u$, we obtain an equation containing on the left-hand side a term $\log\dfrac{\Theta(u - a)}{\Theta(u + a)}$, and on the right-hand an integral in regard to $u$; this, in fact, corresponds to Jacobi's equation
\[
\tfrac{\Theta'a}{\Theta a}\cdot 2u + \log\frac{\Theta(u - a)}{\Theta(u + a)} = -2\int_0^u \frac{k^2\operatorname{sn} a\operatorname{cn} a\operatorname{dn} a\,\operatorname{sn}^2 u\,du}{1 - k^2\operatorname{sn}^2 a\operatorname{sn}^2 u}.
\]

\medskip

\textbf{22.} It may further be noticed that if, in the equation in question and in the three other equations of the system, we introduce into the integral the variable $x$ in place of $u$, and the corresponding quantity $\xi$ in place of $u'$, then the integral is that of an expression such as
\[
\frac{dx}{T\sqrt{a-x\mathbin{.}b-x\mathbin{.}c-x\mathbin{.}d-x}},
\]
where $T$ is $= x - \xi$, or is $= $ any one of three forms such as
\[
\begin{vmatrix} 1 & \xi + x & \xi x \\ 1 & a + b & ab \\ 1 & c + d & cd \end{vmatrix}.
\]

\medskip

\textit{R\'esum\'e of the ulterior theory of the double functions.}

\medskip

\textbf{23.} The ulterior theory of the double functions is intended to be carried out on the like plan. As regards these, it is to be observed here that we have not only the $16$ equations leading to linear relations between the squared functions, but that the remaining $240$ equations lead also to linear relations between binary products of different functions. We have thus between the $16$ functions a system of quadric relations, which in fact determine the ratios of the $16$ functions in terms of two variable parameters $x,y$. (The $16$ functions are thus the coordinates of a point on a quadri-quadric two-fold locus in $15$-dimensional space.) The forms depend upon six constants, $a,b,c,d,e,f$: writing for shortness
\[
\sqrt{a} = \sqrt{a - x\mathbin{.}a - y},
\]
\[
\sqrt{ab} = \tfrac{1}{x - y}\bigl\{\sqrt{a - x\mathbin{.}b - x\mathbin{.}f - x\mathbin{.}c - y\mathbin{.}d - y\mathbin{.}e - y} + \sqrt{a - y\mathbin{.}b - y\mathbin{.}f - y\mathbin{.}c - x\mathbin{.}d - x\mathbin{.}e - x}\bigr\},
\]
(observe that in the symbols $\sqrt{ab}$ it is always $f$ that accompanies the two expressed letters $a,b$ or, what is the same thing, the duad $ab$ is really an abbreviation for the double triad $abf\mathbin{.}cde$): then the $16$ functions are proportional to properly determined constant multiples of
\[
\sqrt{a},\sqrt{b},\sqrt{c},\sqrt{d},\sqrt{e},\sqrt{f},\sqrt{ab},\sqrt{ac},\sqrt{ad},\sqrt{ae},\sqrt{bc},\sqrt{bd},\sqrt{be},\sqrt{cd},\sqrt{ce},\sqrt{de}:
\]
and this suggests that the functions should be represented by the single and double letter notation $A(u,v),\ldots,AB(u,v),\ldots$; viz.\ if for shortness the arguments are omitted, then we have
\[
A, B, C, D, E, F, AB, AC, AD, AE, BC, BD, BE, CD, CE, DE,
\]
proportional to determinate constant multiples of the before-mentioned functions $\sqrt{a},\ldots,\sqrt{ab},\ldots$, of $x$ and $y$.

\medskip

\textbf{24.} It is interesting to notice why in the expressions for $\sqrt{ab}$, \&c., the sign connecting the two radicals is $+$; the effect of the interchange of $x,y$ is, in fact, to change $(u,v)$ into $(-u, -v)$; consequently to change the sign of the odd functions, and to leave unaltered those of the even functions: the interchange does in fact leave $\sqrt{a}$, \&c., unaltered, while it changes $\sqrt{ab}$, \&c., into $-\sqrt{ab}$, \&c.; and thus, since only the ratios are attended to, there is a change of sign as there should be.

\medskip

\textbf{25.} The equations of the product-theorem lead to expressions for
\[
A(u+u')\, B(u-u') - B(u+u')\, A(u-u'),
\]
where the arguments, written above, are used to denote the two arguments, viz.\ $u + u'$ to denote $(u + u', v + v')$ and $u - u'$ to denote $(u - u', v - v')$; and where the letters $A,B$ denote each or either of them a single or double letter. These expressions are found in terms of the functions of $(u,v)$ and of $(u',v')$: in any such expression taking $u',v'$ each of them indefinitely small, but with their ratio arbitrary, we obtain the value of
\[
A(u)\,\partial B(u) - B(u)\,\partial A(u),
\]
(viz.\ $u$ here stands for the two arguments $(u,v)$, and $\partial$ denotes total differentiation
\[
\partial A = du\,\tfrac{d}{du} A(u,v) + dv\,\tfrac{d}{dv} A(u,v)\bigr),
\]
as a quadric function of the functions of $(u,v)$: or dividing by $A^2$, the form is $\partial\dfrac{B}{A}$ equal to a function of the quotient-functions $\dfrac{B}{A}$, \&c., that is, we have the differentials of the quotient-functions in terms of the quotient-functions themselves. Substituting for the quotient-functions their values in terms of $x,y$, we should obtain the differential relations between $dx,dy,du,dv$, viz.\ putting for shortness
\[
X = a - x\mathbin{.}b - x\mathbin{.}c - x\mathbin{.}d - x\mathbin{.}e - x\mathbin{.}f - x,
\]
and
\[
Y = a - y\mathbin{.}b - y\mathbin{.}c - y\mathbin{.}d - y\mathbin{.}e - y\mathbin{.}f - y,
\]
these are of the form
\[
\frac{dx}{\sqrt{X}} - \frac{dy}{\sqrt{Y}},\quad \frac{x\,dx}{\sqrt{X}} - \frac{y\,dy}{\sqrt{Y}},
\]
each of them equal to a linear function of $du$ and $dv$: so that the quotient-functions are in fact the $15$ hyperelliptic functions belonging to the integrals $\int\!\dfrac{dx}{\sqrt{X}},\ \int\!\dfrac{x\,dx}{\sqrt{X}}$, and there is thus an addition-theorem for them, in accordance with the theory of these integrals.

\medskip

\textbf{26.} The first $16$ equations of the product-theorem, putting therein first $u = 0$, $v = 0$, and then $u' = 0, v' = 0$, and using the results to eliminate the functions on the right-hand side, give expressions for
\[
A\mathbin{.}B,\ \&c.,
\]
that is, they give $A(u + u', v + v')\mathbin{.}B(u - u', v - v')$, \&c., in terms of the functions of $(u,v)$ and $(u',v')$: and we have thus an addition-with-subtraction theorem for the double theta-functions. And we have thence also consequences analogous to those which present themselves in the theory of the single functions.

\medskip

\textit{Remark as to notation.}

\medskip

\textbf{27.} I remark, as regards the single theta-functions, that the characteristics
\[
\begin{pmatrix} 0 \\ 0 \end{pmatrix},\ \begin{pmatrix} 1 \\ 0 \end{pmatrix},\ \begin{pmatrix} 0 \\ 1 \end{pmatrix},\ \begin{pmatrix} 1 \\ 1 \end{pmatrix}
\]
might for shortness be represented by a series of current numbers
\[
0, 1, 2, 3:
\]
and the functions be accordingly called $\vartheta_0 u, \vartheta_1 u, \vartheta_2 u, \vartheta_3 u$; but that, instead of this, I prefer to use throughout the before-mentioned functional symbols
\[
A, B, C, D.
\]
As regards the double functions, I do, however, denote the characteristics
\[
\resizebox{\textwidth}{!}{$\displaystyle \begin{pmatrix} 00 \\ 00 \end{pmatrix}\ \begin{pmatrix} 10 \\ 00 \end{pmatrix}\ \begin{pmatrix} 01 \\ 00 \end{pmatrix}\ \begin{pmatrix} 11 \\ 00 \end{pmatrix}\ \Big|\ \begin{pmatrix} 00 \\ 10 \end{pmatrix}\ \begin{pmatrix} 10 \\ 10 \end{pmatrix}\ \begin{pmatrix} 01 \\ 10 \end{pmatrix}\ \begin{pmatrix} 11 \\ 10 \end{pmatrix}\ \Big|\ \begin{pmatrix} 00 \\ 01 \end{pmatrix}\ \begin{pmatrix} 10 \\ 01 \end{pmatrix}\ \begin{pmatrix} 01 \\ 01 \end{pmatrix}\ \begin{pmatrix} 11 \\ 01 \end{pmatrix}\ \Big|\ \begin{pmatrix} 00 \\ 11 \end{pmatrix}\ \begin{pmatrix} 10 \\ 11 \end{pmatrix}\ \begin{pmatrix} 01 \\ 11 \end{pmatrix}\ \begin{pmatrix} 11 \\ 11 \end{pmatrix}$}
\]
by a series of current numbers
\[
0,1,2,3,4,5,6,7,8,9,10,11,12,13,14,15,
\]
and write the functions as $\vartheta_0,\vartheta_1,\ldots,\vartheta_{15}$ accordingly; and I use also, as and when it is convenient, the foregoing single and double letter notation $A, AB,\ldots$, which correspond to them in the order
\[
BD,\ CE,\ CD,\ BE,\ AC,\ C,\ AB,\ B,\ BC,\ DE,\ F,\ A,\ AD,\ D,\ E,\ AE.
\]
Moreover, I write down for the most part a single argument only: thus, $A(u + u')$ stands for $A(u + u', v + v')$, $A(0)$ for $A(0,0)$: and so in other cases.

\bigskip

\begin{center}\textbf{SECOND PART.---THE SINGLE THETA-FUNCTIONS.}\end{center}

\medskip

\textit{Notation, \&c.}

\medskip

\textbf{28.} Writing $\exp.\,\omega = q$, and converting the exponentials into circular functions, we have, directly from the definition,
\begin{align*}
\vartheta\!\begin{pmatrix} 0 \\ 0 \end{pmatrix}\!(u) &= \vartheta_0 u = Au = 1 + 2q\cos\pi u + 2q^4\cos 2\pi u + 2q^9\cos 3\pi u + \ldots,\\
\vartheta\!\begin{pmatrix} 1 \\ 0 \end{pmatrix}\!(u) &= \vartheta_1 u = Bu = 2q^{\tfrac{1}{4}}\cos\tfrac{1}{2}\pi u + 2q^{\tfrac{9}{4}}\cos\tfrac{3}{2}\pi u + 2q^{\tfrac{25}{4}}\cos\tfrac{5}{2}\pi u + \ldots,\\
\vartheta\!\begin{pmatrix} 0 \\ 1 \end{pmatrix}\!(u) &= \vartheta_2 u = Cu = 1 - 2q\cos\pi u + 2q^4\cos 2\pi u - 2q^9\cos 3\pi u + \ldots\ (= \Theta(Ku),\ \text{Jacobi}),\\
\vartheta\!\begin{pmatrix} 1 \\ 1 \end{pmatrix}\!(u) &= \vartheta_3 u = Du = -2q^{\tfrac{1}{4}}\sin\tfrac{1}{2}\pi u + 2q^{\tfrac{9}{4}}\sin\tfrac{3}{2}\pi u - 2q^{\tfrac{25}{4}}\sin\tfrac{5}{2}\pi u + \ldots\ (= -H(Ku),\ \text{Jacobi}),
\end{align*}
where $\omega$ is of the form $\omega = -\alpha + \beta i$, $\alpha$ being non-evanescent and positive: hence $q = \exp.\,(-\alpha + \beta i) = e^{-\alpha}(\cos\beta + i\sin\beta)$, where $e^{-\alpha}$, the modulus of $q$, is positive and less than $1$; $\cos\beta$ may be either positive or negative, and $q^{\tfrac{1}{4}}$ is written to denote $\exp.\,\tfrac{1}{4}(-\alpha + \beta i)$, viz.\ this is $= e^{-\tfrac{1}{4}\alpha}(\cos\tfrac{1}{4}\beta + i\sin\tfrac{1}{4}\beta)$. But usually $\beta = 0$, viz.\ $q$ is real positive quantity less than $1$, and $q^{\tfrac{1}{4}}$ denotes the real fourth root of $q$.

\medskip

I have given above the three notations but, as already mentioned, I propose to employ for the four functions the notation $Au, Bu, Cu, Du$: it will be observed that $Du$ is an odd function, but that $Au, Bu, Cu$ are even functions, of $u$.

\medskip

\textit{The constants of the theory.}

\medskip

\textbf{29.} We have
\begin{align*}
A0 &= 1 + 2q + 2q^4 + 2q^9 + \ldots,\\
B0 &= 2q^{\tfrac{1}{4}} + 2q^{\tfrac{9}{4}} + 2q^{\tfrac{25}{4}} + \ldots,\\
C0 &= 1 - 2q + 2q^4 - 2q^9 + \ldots,\\
D0 &= 0,\\
D'0 &= -\pi\{q^{\tfrac{1}{4}} - 3q^{\tfrac{9}{4}} + 5q^{\tfrac{25}{4}} - \ldots\}.
\end{align*}
If, as definitions of $k, k', K$, we assume
\[
k = \frac{B^2 0}{A^2 0},\quad k' = \frac{C^2 0}{A^2 0},\quad K = \frac{A^2 0}{B0\, C0}\cdot \frac{D'0}{C0},
\]
then we have
\[
k = 4\sqrt{q}\,\Bigl\{\frac{1 + q^4 + q^8 + \ldots}{1 + 2q + 2q^4 + \ldots}\Bigr\}^2,\ \&c.\ \ldots = 4\sqrt{q}\,(1 - 4q + 14q^2 + \ldots),
\]
where I have added the first few terms of the expansions of these quantities. We have identically
\[
k^2 + k'^2 = 1.
\]
It will be convenient to write also, as the definition of $E$,
\[
\frac{E}{K} = \frac{1}{K^2}\cdot \frac{C''0}{C0} = \frac{2\pi^2}{K^2}\cdot \frac{q - 4q^4 + 9q^9 - \ldots}{1 - 2q + 2q^4 - \ldots};
\]
we have then
\[
\frac{E}{K} = 1 - 8q + 48q^2 - 224q^3 + \ldots,
\]
moreover,
\[
E = \tfrac{1}{2}\pi\bigl\{1 - 4q + 20q^2 - 64q^3 + \ldots\bigr\}.
\]
giving

\medskip

\textbf{30.} Other formulae are
\[
k = 4\sqrt{q}\,\Bigl(\frac{1 + q^2\mathbin{.}1 + q^4 \ldots}{1 + q\mathbin{.}1 + q^3 \ldots}\Bigr)^4,
\]
\[
k' = \Bigl(\frac{1 - q\mathbin{.}1 - q^3 \ldots}{1 + q\mathbin{.}1 + q^3 \ldots}\Bigr)^4,
\]
\[
K = \tfrac{1}{2}\pi\bigl(1 + q\mathbin{.}1 + q^3\mathbin{.}1 + q^5 \ldots\bigr)^4.
\]

\medskip

\textbf{31.} Jacobi's definition of $q$ is from a different point of view altogether, viz.\ we have $q = \exp.\,-\dfrac{\pi K'}{K}$, where
\[
K = \int_0^{\tfrac{1}{2}\pi}\!\frac{d\phi}{\sqrt{1 - k^2\sin^2\phi}},
\]
and $K'$ is the like function of $k'$; the equation gives $\log q = -\dfrac{\pi K'}{K}$, viz.\ we have
\[
\frac{K'}{K} = -\dfrac{1}{\pi}\log q,
\]
and, regarding herein $k$ as a given function of $q$, this equation gives $\dfrac{K'}{K}$ as a function of $q$.

\medskip

\textit{The product-theorem.}

\medskip

\textbf{32.} The product-theorem is
\[
\resizebox{\textwidth}{!}{$\displaystyle \vartheta\!\begin{pmatrix} \alpha \\ \gamma \end{pmatrix}\!(u + u')\,\vartheta\!\begin{pmatrix} \alpha' \\ \gamma' \end{pmatrix}\!(u - u') = \Theta\!\begin{pmatrix} \tfrac{1}{2}(\alpha + \alpha') \\ \gamma + \gamma' \end{pmatrix}\!(2u)\cdot \Theta\!\begin{pmatrix} \tfrac{1}{2}(\alpha - \alpha') \\ \gamma - \gamma' \end{pmatrix}\!(2u') + \Theta\!\begin{pmatrix} \tfrac{1}{2}(\alpha + \alpha') + 1 \\ \gamma + \gamma' \end{pmatrix}\!(2u)\cdot \Theta\!\begin{pmatrix} \tfrac{1}{2}(\alpha - \alpha') + 1 \\ \gamma - \gamma' \end{pmatrix}\!(2u')$}.
\]
Here giving to $\alpha,\alpha',\gamma,\gamma'$ their different values, and introducing unaccented and accented capitals to denote the functions of $2u$ and $2u'$ respectively, the $16$ equations are

\medskip

\textit{Square-set:}
\begin{align*}
A\mathbin{.}A\quad &\vartheta\!\begin{pmatrix} 0 \\ 0 \end{pmatrix}\!(u + u')\,\vartheta\!\begin{pmatrix} 0 \\ 0 \end{pmatrix}\!(u - u') = XX' + YY',\\
B\mathbin{.}B\quad &\vartheta\!\begin{pmatrix} 1 \\ 0 \end{pmatrix}\,,,\,\vartheta\!\begin{pmatrix} 1 \\ 0 \end{pmatrix}\,, = YX' + XY',\\
C\mathbin{.}C\quad &\vartheta\!\begin{pmatrix} 0 \\ 1 \end{pmatrix}\,,,\,\vartheta\!\begin{pmatrix} 0 \\ 1 \end{pmatrix}\,, = XX' - YY',\\
D\mathbin{.}D\quad &\vartheta\!\begin{pmatrix} 1 \\ 1 \end{pmatrix}\,,,\,\vartheta\!\begin{pmatrix} 1 \\ 1 \end{pmatrix}\,, = -YX' + XY';
\end{align*}

\textit{First product-set:}
\begin{align*}
C\mathbin{.}A\quad &\vartheta\!\begin{pmatrix} 0 \\ 1 \end{pmatrix}\!(u + u')\,\vartheta\!\begin{pmatrix} 0 \\ 0 \end{pmatrix}\!(u - u') = X_1 X_1' + Y_1 Y_1',\\
A\mathbin{.}C\quad &\vartheta\!\begin{pmatrix} 0 \\ 0 \end{pmatrix}\,,,\,\vartheta\!\begin{pmatrix} 0 \\ 1 \end{pmatrix}\,, = X_1 X_1' - Y_1 Y_1',\\
D\mathbin{.}B\quad &\vartheta\!\begin{pmatrix} 1 \\ 1 \end{pmatrix}\,,,\,\vartheta\!\begin{pmatrix} 1 \\ 0 \end{pmatrix}\,, = Y_1 X_1' + X_1 Y_1',\\
B\mathbin{.}D\quad &\vartheta\!\begin{pmatrix} 1 \\ 0 \end{pmatrix}\,,,\,\vartheta\!\begin{pmatrix} 1 \\ 1 \end{pmatrix}\,, = Y_1 X_1' - X_1 Y_1';
\end{align*}

\textit{Second product-set:}
\begin{align*}
B\mathbin{.}A\quad &\vartheta\!\begin{pmatrix} 1 \\ 0 \end{pmatrix}\!(u + u')\,\vartheta\!\begin{pmatrix} 0 \\ 0 \end{pmatrix}\!(u - u') = PP' + QQ',\\
A\mathbin{.}B\quad &\vartheta\!\begin{pmatrix} 0 \\ 0 \end{pmatrix}\,,,\,\vartheta\!\begin{pmatrix} 1 \\ 0 \end{pmatrix}\,, = PQ' + QP',\\
D\mathbin{.}C\quad &\vartheta\!\begin{pmatrix} 1 \\ 1 \end{pmatrix}\,,,\,\vartheta\!\begin{pmatrix} 0 \\ 1 \end{pmatrix}\,, = iPP' - iQQ',\\
C\mathbin{.}D\quad &\vartheta\!\begin{pmatrix} 0 \\ 1 \end{pmatrix}\,,,\,\vartheta\!\begin{pmatrix} 1 \\ 1 \end{pmatrix}\,, = iPQ' - iQP';
\end{align*}

\textit{Third product-set:}
\begin{align*}
D\mathbin{.}A\quad &\vartheta\!\begin{pmatrix} 1 \\ 1 \end{pmatrix}\!(u + u')\,\vartheta\!\begin{pmatrix} 0 \\ 0 \end{pmatrix}\!(u - u') = P_1 P_1' + Q_1 Q_1',\\
A\mathbin{.}D\quad &\vartheta\!\begin{pmatrix} 0 \\ 0 \end{pmatrix}\,,,\,\vartheta\!\begin{pmatrix} 1 \\ 1 \end{pmatrix}\,, = iP_1 Q_1' - iQ_1 P_1',\\
B\mathbin{.}C\quad &\vartheta\!\begin{pmatrix} 1 \\ 0 \end{pmatrix}\,,,\,\vartheta\!\begin{pmatrix} 0 \\ 1 \end{pmatrix}\,, = -iP_1 P_1' + iQ_1 Q_1',\\
C\mathbin{.}B\quad &\vartheta\!\begin{pmatrix} 0 \\ 1 \end{pmatrix}\,,,\,\vartheta\!\begin{pmatrix} 1 \\ 0 \end{pmatrix}\,, = P_1 Q_1' + Q_1 P_1'.
\end{align*}

\medskip

\textbf{33.} Here, and subsequently, we have
\[
\Theta\!\begin{pmatrix} 0 \\ 0 \end{pmatrix}, \Theta\!\begin{pmatrix} 1 \\ 0 \end{pmatrix}, \Theta\!\begin{pmatrix} 0 \\ 1 \end{pmatrix}, \Theta\!\begin{pmatrix} 1 \\ 1 \end{pmatrix}\!(2u) = X, Y, X_1, Y_1\quad \mid\quad \Theta\!\begin{pmatrix} \tfrac{1}{2} \\ 0 \end{pmatrix}, \Theta\!\begin{pmatrix} \tfrac{3}{2} \\ 0 \end{pmatrix}, \Theta\!\begin{pmatrix} \tfrac{1}{2} \\ 1 \end{pmatrix}, \Theta\!\begin{pmatrix} \tfrac{3}{2} \\ 1 \end{pmatrix}\!(2u) = P, Q, P_1, Q_1,
\]
\[
\Theta\ldots(2u') = X', Y', X_1', Y_1';\quad \Theta\ldots(2u') = P', Q', P_1', Q_1',
\]
\[
\Theta\ldots(0) = \alpha, \beta, \alpha_1, \beta_1;\quad \Theta\ldots(0) = p, q, p_1, q_1;
\]
viz.\ we use also $\alpha,\beta,\alpha_1,\beta_1$ and $p,q,p_1,q_1$ to denote the zero-functions; $\beta_1$ is $= 0$, but we use $\beta_1'$ to denote the zero-value of $\dfrac{d}{du}Y_1$.

\medskip

\textbf{34.} In order to obtain the foregoing relations, it is necessary to observe that
\[
\Theta\!\begin{pmatrix} \alpha + 2 \\ \gamma \end{pmatrix},,\, = \Theta\!\begin{pmatrix} \alpha \\ \gamma \end{pmatrix},,\,;
\]
by which the upper character is always reduced to $0, 1, \tfrac{1}{2}$, or $\tfrac{3}{2}$; and that, for reducing the lower character, we have
\[
\Theta\!\begin{pmatrix} \alpha \\ \gamma + 2 \end{pmatrix},, = \Theta\!\begin{pmatrix} \alpha \\ \gamma \end{pmatrix},,;\quad \Theta\!\begin{pmatrix} \tfrac{1}{2} \\ \gamma + 2 \end{pmatrix},, = -\Theta\!\begin{pmatrix} \tfrac{1}{2} \\ \gamma \end{pmatrix},,;
\]
\[
\Theta\!\begin{pmatrix} 1 \\ \gamma + 2 \end{pmatrix},, = i\Theta\!\begin{pmatrix} 1 \\ \gamma \end{pmatrix},,;\quad \Theta\!\begin{pmatrix} \tfrac{3}{2} \\ \gamma - 2 \end{pmatrix},, = -i\Theta\!\begin{pmatrix} \tfrac{3}{2} \\ \gamma \end{pmatrix},;\quad \Theta\!\begin{pmatrix} \tfrac{3}{2} \\ \gamma + 2 \end{pmatrix},, = \Theta\!\begin{pmatrix} \tfrac{3}{2} \\ \gamma \end{pmatrix},;
\]
by means of which the lower character is always reduced to $0$ or $1$: in all these formulae the argument is arbitrary, and it is thus $= 2u$, or $2u'$ as the case requires. The formulae are obtained without difficulty directly from the definition of the functions $\Theta$.

\medskip

\textbf{35.} As an instance, taking $\dfrac{\alpha}{\gamma},\dfrac{\alpha'}{\gamma'} = \dfrac{1}{1},\dfrac{0}{1}$, the product-equation is
\[
\vartheta\!\begin{pmatrix} 1 \\ 1 \end{pmatrix}\!(u + u')\,\vartheta\!\begin{pmatrix} 0 \\ 1 \end{pmatrix}\!(u - u') = \Theta\!\begin{pmatrix} \tfrac{1}{2} \\ 2 \end{pmatrix}\!(2u)\cdot \Theta\!\begin{pmatrix} \tfrac{1}{2} \\ 0 \end{pmatrix}\!(2u') + \Theta\!\begin{pmatrix} \tfrac{3}{2} \\ 2 \end{pmatrix}\!(2u)\cdot \Theta\!\begin{pmatrix} \tfrac{3}{2} \\ 0 \end{pmatrix}\!(2u'),
\]
\[
= -i\Theta\!\begin{pmatrix} \tfrac{1}{2} \\ 0 \end{pmatrix}\!(2u)\cdot \Theta\!\begin{pmatrix} \tfrac{1}{2} \\ 0 \end{pmatrix}\!(2u') - i\Theta\!\begin{pmatrix} \tfrac{3}{2} \\ 0 \end{pmatrix}\!(2u)\cdot \Theta\!\begin{pmatrix} \tfrac{3}{2} \\ 0 \end{pmatrix}\!(2u'),
\]
\[
= iPP' - iQQ',
\]
which agrees with the before-given value.

\medskip

\textbf{36.} The following values are not actually required: but I give them to fix the ideas and to show the meaning of the quantities with which we work.

\medskip

\noindent\resizebox{\textwidth}{!}{\begin{tabular}{l@{\hspace{1em}}l}
$X = \Theta\!\begin{pmatrix} 0 \\ 0 \end{pmatrix}\!(2u) = 1 + 2q^4\cos 2\pi u + 2q^{16}\cos 4\pi u + \ldots,$ & $u = 0:$\ \ \ $\alpha = 1 + 2q^4 + 2q^{16} + \ldots,$\\[4pt]
$Y = \Theta\!\begin{pmatrix} 1 \\ 0 \end{pmatrix}\!(2u) = 2q\cos\pi u + 2q^9\cos 3\pi u + \ldots,$ & $\beta = 2q + 2q^9 + \ldots,$\\[4pt]
$X_1 = \Theta\!\begin{pmatrix} 0 \\ 1 \end{pmatrix}\!(2u) = 1 - 2q^4\cos 2\pi u + 2q^{16}\cos 4\pi u - \ldots,$ & $\alpha_1 = 1 - 2q^4 + 2q^{16} - \ldots,$\\[4pt]
$Y_1 = \Theta\!\begin{pmatrix} 1 \\ 1 \end{pmatrix}\!(2u) = -2q\sin\pi u + 2q^9\sin 3\pi u - \ldots,$ & $\beta_1' = 2\pi\,(-q + 3q^9 - \ldots) = \dfrac{d}{du}Y_1\ \text{for}\ u = 0.$
\end{tabular}}

\medskip

\begin{align*}
P &= \Theta\!\begin{pmatrix} \tfrac{1}{2} \\ 0 \end{pmatrix}\!(2u) = q^{\tfrac{1}{4}}(\cos\tfrac{1}{2}\pi u + i\sin\tfrac{1}{2}\pi u) + q^{\tfrac{9}{4}}(\cos\tfrac{3}{2}\pi u - i\sin\tfrac{3}{2}\pi u)\\
&\hspace{6em}+ q^{\tfrac{25}{4}}(\cos\tfrac{5}{2}\pi u + i\sin\tfrac{5}{2}\pi u) + \ldots,\\
Q &= \Theta\!\begin{pmatrix} \tfrac{3}{2} \\ 0 \end{pmatrix}\!(2u) = q^{\tfrac{1}{4}}(\cos\tfrac{1}{2}\pi u - i\sin\tfrac{1}{2}\pi u) + q^{\tfrac{9}{4}}(\cos\tfrac{3}{2}\pi u + i\sin\tfrac{3}{2}\pi u)\\
&\hspace{6em}+ q^{\tfrac{25}{4}}(\cos\tfrac{5}{2}\pi u - i\sin\tfrac{5}{2}\pi u) + \ldots,\\
P_1 &= \Theta\!\begin{pmatrix} \tfrac{1}{2} \\ 1 \end{pmatrix}\!(2u) = \tfrac{1 + i}{\sqrt{2}}\bigl\{q^{\tfrac{1}{4}}(\cos\tfrac{1}{2}\pi u + i\sin\tfrac{1}{2}\pi u) - q^{\tfrac{9}{4}}(\cos\tfrac{3}{2}\pi u - i\sin\tfrac{3}{2}\pi u)\\
&\hspace{8em}- q^{\tfrac{25}{4}}(\cos\tfrac{5}{2}\pi u + i\sin\tfrac{5}{2}\pi u) + \ldots\bigr\},\\
Q_1 &= \Theta\!\begin{pmatrix} \tfrac{3}{2} \\ 1 \end{pmatrix}\!(2u) = \tfrac{1 - i}{\sqrt{2}}\bigl\{q^{\tfrac{1}{4}}(\cos\tfrac{1}{2}\pi u - i\sin\tfrac{1}{2}\pi u) - q^{\tfrac{9}{4}}(\cos\tfrac{3}{2}\pi u + i\sin\tfrac{3}{2}\pi u)\\
&\hspace{8em}- q^{\tfrac{25}{4}}(\cos\tfrac{5}{2}\pi u - i\sin\tfrac{5}{2}\pi u) + \ldots\bigr\};
\end{align*}
and therefore also
\[
p = q = q^{\tfrac{1}{4}} + q^{\tfrac{9}{4}} + q^{\tfrac{25}{4}} + \ldots,
\]
\[
p_1 = \tfrac{1 + i}{\sqrt{2}}\bigl\{q^{\tfrac{1}{4}} - q^{\tfrac{9}{4}} - q^{\tfrac{25}{4}} + q^{\tfrac{49}{4}} + q^{\tfrac{81}{4}} - \ldots\bigr\},\quad q_1 = \tfrac{1 - i}{\sqrt{2}}\bigl\{\text{Do.}\bigr\};\ p_1 = iq_1.
\]

\end{document}
